\begin{figure}[!hbp]
\begin{examplepanel}{tgtcol}{Factory and certificates in simplified target notation}
\small
\[
\tgt{\begin{aligned}
 v_{\mathrm{in}}
   &=\lambda^{\{\ev{console}\}}(v:\ty{Unit}).\
      \ev{console}\ v\ \langle g_{\mathrm{in}}\rangle,\\
 v_{\mathrm{make}}
   &=\lambda^{\{\ev{log}\}}(u:\ty{Unit}).\\[-2pt]
   &\quad
      \bigl(\ev{let}\ r=\ev{log}\ u\ \ev{in}\
         v_{\mathrm{in}}\ \langle\varnothing;q_0\rangle\bigr)
       \langle g_{\mathrm{make}}\rangle .
\end{aligned}}
\]
The symbol \(v_{\mathrm{in}}\) abbreviates the lambda literal itself,
with empty immediate uses. Its certificates and those of the factory are:
\[
\tgt{\begin{aligned}
 g_{\mathrm{in}}&=\ev{elem}(\{\ev{console},v\},\{\ev{console},v\}),\\
 q_0&=\ev{elem}(\varnothing,\varnothing),\\
 g_{\mathrm{make}}&=\ev{elem}(\{\ev{log},u\},\{\ev{log},u\}).
\end{aligned}}
\]
In a let annotation \(\langle U;q\rangle\), \(U\) bounds the continuation
body's uses and \(q\) proves the avoidance obligation.
The named certificates are checked under their displayed binders:
\(g_{\mathrm{in}}\) under \(v\), \(g_{\mathrm{make}}\) under \(u\),
and \(q_0\) under \(r\).
\end{examplepanel}
\begin{examplepanel}{tgtcol}{Clients of the stored factory \(\ev{make}\)}
\small
\[
\tgt{\begin{aligned}
 P_{\mathrm{call}}
   &=\ev{let}\ f=\ev{make}\ \ev{unit}\ \ev{in}\
      f\ \ev{unit}\ \langle U;q_{\mathrm{call}}\rangle,\\
 P_{\mathrm{drop}}
   &=\ev{let}\ f=\ev{make}\ \ev{unit}\ \ev{in}\
      \ev{unit}\ \langle\varnothing;q_{\mathrm{drop}}\rangle,\\
 U&=\{\ev{console},\ev{log}\},\\
 q_{\mathrm{call}}&:\{f,\ev{unit}\}\le_c U,\qquad
 q_{\mathrm{drop}}=\ev{capvar}(\ev{unit}):\{\ev{unit}\}\le_c\varnothing.
\end{aligned}}
\]
\end{examplepanel}
\begin{examplepanel}{runcol}{Initial roots of the whole state}
\[
\run{P_{\mathrm{call}}:\{\kappa_\ell,\kappa_c\},
      \qquad P_{\mathrm{drop}}:\{\kappa_\ell\}.}
\]
\end{examplepanel}
\caption{A factory with distinct immediate uses and result captures.
Both clients call \(\mathsf{make}\), which invokes \(\mathsf{log}\) and
returns a closure that invokes \(\mathsf{console}\);
\(P_{\mathrm{call}}\) then invokes that result.
The text constructs \(q_{\mathrm{call}}\) from two branches. This named
presentation follows the C1/C6 examples, with certificates supplied
directly as target evidence. Source translation may produce different
evidence syntax.}
\label{fig:factory}
\end{figure}
